\documentclass{article}
\usepackage[T1]{fontenc}
\usepackage{lmodern}
\usepackage{graphicx}
\usepackage{amsmath,amssymb}
\usepackage[a4paper,margin=2cm]{geometry}
\usepackage[absolute,overlay]{textpos}
\setlength{\TPHorizModule}{1mm}
\setlength{\TPVertModule}{1mm}
\usepackage[indonesian]{babel}
\usepackage{hyperref}
\setlength{\emergencystretch}{2em}
% unit: Halaman 37

\begin{document}

% === Halaman 37 (python/native) ===

37

• Hasil enkripsi seluruhnya adalah sebagai berikut:

 Plainteks 

: terimapesanangulaikambing

 Kunci 

: soreangsoreangsoreangsore

 Cipherteks : LSIMMNVWGRRAAMMZRMKNSTWEK

• Tanpa menggunakan Vigenere Square pun enkripsi tetap dapat dihitung secara Caesar 

Cipher dengan menjumlahkan  plainteks pj  dengan kunci ki dalam modulus 26:

              Enkripsi:  cj = E(pj) = (pj + ki) mod 26   

(1)


Dekripsi:  pj = D(cj) = (cj – ki) mod 26   

(2)

    Contoh:


(t + s)  mod 26 = (19 + 18) mod 26 = 37 mod 26 = 11 =  L


(e + o) mod 26 = (4 + 14) mod 26 = 18 mod 26 = 18 = S, dst

A = 0,  B = 1, C = 2,  D = 3,  E = 4,  F = 5,  G = 6,  H = 7,  I = 8,  J = 9,  K = 10

L = 11, M = 12, N = 13, O = 14, P = 15, Q = 16, R = 17, S = 18, T = 19, U = 20

V = 21, W = 22, X = 23, Y = 24, Z = 25

\end{document}
